\section{Introduction}
\label{sec:intro}

\IEEEPARstart{C}{ompact} models, the equations a circuit simulator evaluates for every
transistor, diode or photonic element at every Newton iteration, are distributed as Verilog-A
source~\cite{accellera-vams-lrm}. Industrial models such as PSP~\cite{gildenblat2006psp} and
BSIM are standardized by the Compact Model Coalition and released in this form. Running one
requires a compiler that turns the source into code a simulator can evaluate, together with the
partial derivatives Newton's method needs.

\todo{Motivation paragraph: why another Verilog-A tool. Candidate angles: (i) a permissively
licensed (Apache-2.0), clean-room implementation that other projects can embed; (ii) no native
toolchain, so a single static binary with nothing to install; (iii) bit-for-bit reproducibility
across operating systems, CPUs and thread counts, which general-purpose simulators do not claim;
(iv) multi-physics through Verilog-A disciplines (thermal, optical, mechanical) in one system.
Pick the one or two the paper leads with; (iii) is the most novel measured result.}

This paper makes the following contributions:
\begin{itemize}
  \item An architecture that splits a Verilog-A simulator along two frozen interfaces, so that
  the numerical core is developed and validated against hand-written models independently of the
  language compiler (Section~\ref{sec:arch}).
  \item A forward-mode automatic-differentiation back end whose every operator is checked against
  central finite differences, together with the numerical pitfalls met on production compact
  models (Section~\ref{sec:codegen}).
  \item A DC rescue sequence (plain Newton, adaptive $g_\text{min}$ stepping, damping, per-node
  step capping) and a block-triangular sparse solve that bring transistor-level ISCAS'85
  circuits of up to \num{268676} unknowns to verified operating points (Sections~\ref{sec:dc}
  and~\ref{sec:scale}).
  \item A set of design rules that make the output bit-identical across Linux, macOS and Windows,
  x86-64 with and without AVX-512 and Apple arm64, and any thread count, with measurements that
  show it holds (Section~\ref{sec:determinism}).
  \item A validation methodology in which every analysis is gated against committed QSPICE
  reference outputs, and results at scale are checked against an independent oracle, the Boolean
  function of the circuit (Sections~\ref{sec:validation} and~\ref{sec:scale}).
  \item A case study showing that Verilog-A written for an OSDI flow runs unchanged and
  standalone: two electro-optic micro-ring modulator models agree with OpenVAF and ngspice on the
  same source, and one reproduces the main results of the published model it implements, without
  the solver-specific rewrites ngspice needed (Section~\ref{sec:mrm}).
\end{itemize}

\todo{Roadmap paragraph for the rest of the paper, once the section order is fixed.}
